% ID: SEL-D0012-BA-E02
% ESTADO: cotejada. Segunda pasada visual completa y borrador compilado revisados.
% Fuente: Selectividad/matematicas_aplicadas/EXTRA_RESERVA_B_Modelo-6-Enunciados (Mat-Ap-2021).pdf
% SHA-256: 1752ec2f46b4ef2cc5d1c3346d69b79c2cc8a7aa12f6a0cd8f0d994b5bcda157
% Página: 1; bloque A; ejercicio original 2.

Se consideran las matrices
$A=\begin{pmatrix}1&0&1\\0&1&0\\1&0&1\end{pmatrix}$,
$B=\begin{pmatrix}1&0&2\\1&1&-1\\2&1&0\end{pmatrix}$ y
$C=\begin{pmatrix}1\\-3\\1\end{pmatrix}$.
\begin{enumerate}[label=\alph*)]
  \item \textbf{(1 punto)} Calcule $A^2$, $A^3$, $A^4$ y deduzca la
  expresión de $A^n$, con $n$ un número natural.
  \item \textbf{(0.5 puntos)} Razone si existe la inversa de la matriz $B$.
  \item \textbf{(1 punto)} Razone si la ecuación matricial $B\cdot X=C$
  tiene solución y resuélvala en caso de que sea posible.
\end{enumerate}
